% English version of paper/ja_v4/07_結論.md (ver4, 2026-11-17): one paragraph, no numbers repeated from the abstract (MIT CommLab, Widom), ending with how the pseudo-labels are built; last sentence made positive on 2026-11-17 night (researcher's comment), with the comparator stated (reviews/v4_ja/numbers_intro_rq_2026-11-17.md, item 12).
\section{Conclusion}\label{sec:conclusion}

For multimodal emotion recognition on a new corpus with only a few labels per emotion, we compared distillation on pseudo-labels including other modalities with domain adaptation, semi-supervised and test-time adaptation methods that do not use such pseudo-labels, under the same conditions, with linear classifiers on frozen encoders. For audio classifiers, retraining on pseudo-labels including other modalities gave higher UAR than the compared methods, and, across pseudo-labels from different providers, retrained classifiers had higher UAR the more accurate the pseudo-labels were than the classifier before retraining. In the all-source setting, same-modality unlabelled utterances rarely reduced the remaining error (the error left after correcting prediction bias). In addition, to make the pseudo-labels more accurate, we proposed late fusion that sets the weights of each modality's predictions with strengths from the few target labels and per-emotion profiles from the source corpora. The proposal produced more accurate pseudo-labels than equal-weight fusion. On a new corpus with few labels, the other modalities of the same utterance made up for part of the missing labels, and weighting the modalities with the few target labels and the source corpora, rather than equally, gave AVT classifiers retrained on the resulting pseudo-labels higher UAR in most conditions.
